\documentclass[11pt]{article}
\usepackage{preamble}
\renewcommand{\answer}[1]{}

\begin{document}

\ladderheading{AP Statistics \textperiodcentered\ Unit 4 \textperiodcentered\ Topic 4.7 \textperiodcentered\ B: 2027-01-25 \textperiodcentered\ E: 2027-03-12}{Two Samples of Service Times}

\focusbanner{TODAY'S PROBLEM}

Suppose two depots independently take SRSs of 36 weekday service-time records from monthly populations of 900 North and 1,200 South records. Technology gives $df=69.31$ and $t^*=1.995$ for a 95\% two-sample interval.\par\smallskip \textbf{Tariq:} ``I sorted both samples and paired equal ranks. The resulting differences have $\bar d=2.2$ and $s_d=3.0$, giving $(1.18,3.22)$. Pairing makes the estimate more precise.''\par \textbf{Mei:} ``The records were sampled separately, so I would use a two-sample interval. Sorting observed responses may change the uncertainty rather than reveal design-based pairs.''\par
\begin{center}
\textbf{Service-time samples (min)}\par\smallskip
\begin{tabular}{|l|c|c|c|c|}
\hline
\textbf{Depot} & \textbf{$N$} & \textbf{$n$} & \textbf{$\bar x$} & \textbf{$s$} \\ \hline
\textbf{North} & 900 & 36 & 42.3 & 8.4 \\ \hline
\textbf{South} & 1,200 & 36 & 40.1 & 7.6 \\ \hline
\end{tabular}
\end{center}
\begin{firsttakebox}{FIRST TAKE --- before the video (5 min)}
Does sorting the times make the North and South records real pairs? Explain in one sentence.
\workspace{0.9in}
\end{firsttakebox}

\begin{callout}{\IconBook\ LET THE DESIGN CHOOSE THE INTERVAL (keep on the page)}{calloutgreen}
For two independent random samples, use a two-sample $t$ interval:
$(\bar x_1-\bar x_2)\pm t^*\sqrt{s_1^2/n_1+s_2^2/n_2}$. Paired measurements have a link
in the study design, such as two measurements from the same person; they use one sample of differences.
Check independent random samples or assignment, each 10 percent condition for sampling without replacement,
and both sample shapes or both sample sizes over 30. Interpret in the stated subtraction order.
An interval containing 0 leaves equal population means plausible; an interval entirely on one side supports a difference.
\end{callout}

\newpage
\textbf{Finish after the video:}\par\smallskip

\dokbadge{1}~\textbf{(a)} Find the North-minus-South difference in sample mean service times.\par
\workspace{0.8in}

\dokbadge{2}~\textbf{(b)} The independent-sample standard error is 1.88797 minutes. Use $t^*=1.995$ to calculate the 95\% interval.\par
\workspace{1.4in}

\dokbadge{3}~\textbf{(c)} In 3--4 sentences, choose a procedure using the original sampling design. Compare the two intervals and explain how sorting the times could mislead a reader about whether the population means differ.\par
\workspace{2.0in}

\begin{sentenceframebox}\raggedright\textbf{Frame:} The design supports a \blank[1.6in] interval because \blank[2.2in]; sorting ranks instead \blank[2.0in].\end{sentenceframebox}

\begin{sentenceframebox}\raggedright\textbf{Frame:} The two uncertainty estimates are \blank[1.6in] and \blank[1.6in], which changes \blank[2.0in].\end{sentenceframebox}

\begin{summaryexitbox}{EXIT --- turn this in}
One thing the video changed about my first take: \hrulefill\par
\workspace{0.4in}
\end{summaryexitbox}

\end{document}
